\subsection{Comparative Evaluation Against SOTA Baselines on Two Datasets}
\label{sec:res_h}

To assess the proposed \textit{MicroPhys-BMS} against SOTA paradigms, a comparative evaluation was conducted on two LFP datasets: Stanford ($25^\circ$C, 15{,}100 samples) \cite{che2025enhanced} and CALCE A123 ($24^\circ$C, 11{,}968 samples) \cite{calce2014a123, boyce2022calce}. Three benchmarks were implemented under identical protocols (70/30 split, same 9-D features): (i) PINN with 4 layers of 64 Tanh neurons \cite{karniadakis2021physics}; (ii) Tabular Transformer ($d_{\text{model}} = 27$, $n_{\text{head}} = 3$) \cite{chen2024transformer}; (iii) 5-model LightGBM ensemble (2{,}500 trees total) \cite{ke2017lightgbm}.

\subsubsection{Cross-Dataset Accuracy Comparison}
Table~\ref{tab:comparison_results} reports performance. On Stanford LFP, MicroPhys-BMS (961 params) achieves MAE $= 0.536\%$, competitive with LightGBM ($0.394\%$, $157{,}500$ params, $164\times$ more). PINN achieves $0.898\%$ ($13{,}185$ params), Transformer $1.188\%$ ($13{,}899$ params). On CALCE A123 (independent manufacturer, 18650 form), the proposed model achieves MAE $= 0.173\%$ with $98.89\%$ CDF within $|\Delta\text{SOC}| \le 2.0\%$, matching neural baselines with $13.7\times$--$164\times$ fewer parameters.

\begin{table}[!t]
	\caption{Cross-Dataset Comparative Evaluation}
	\label{tab:comparison_results}
	\centering
	\resizebox{\columnwidth}{!}{%
		\begin{tabular}{lcccc}
			\toprule
			\textbf{Model} & \textbf{Params} & \textbf{Dataset} & \textbf{MAE} & \textbf{RMSE} \\
			\midrule
			\multirow{2}{*}{MicroPhys-BMS (Ours)} & \multirow{2}{*}{\textbf{961}} & Stanford & $0.536\%$ & $1.492\%$ \\
			 & & CALCE A123 & $\mathbf{0.173\%}$ & $2.451\%$ \\
			\midrule
			\multirow{2}{*}{PINN \cite{karniadakis2021physics}} & \multirow{2}{*}{$13{,}185$} & Stanford & $0.898\%$ & $2.050\%$ \\
			 & & CALCE A123 & $0.137\%$ & $2.477\%$ \\
			\midrule
			\multirow{2}{*}{Transformer \cite{chen2024transformer}} & \multirow{2}{*}{$13{,}899$} & Stanford & $1.188\%$ & $2.367\%$ \\
			 & & CALCE A123 & $0.132\%$ & $2.484\%$ \\
			\midrule
			\multirow{2}{*}{LightGBM Ens. \cite{ke2017lightgbm}} & \multirow{2}{*}{$\approx 157{,}500$} & Stanford & $\mathbf{0.394\%}$ & $\mathbf{1.037\%}$ \\
			 & & CALCE A123 & $0.231\%$ & $1.299\%$ \\
			\bottomrule
		\end{tabular}%
	}
\end{table}

\subsubsection{Noise Robustness and CDF Analysis}
As depicted in Fig.~\ref{fig:cdf_comparison}(a), MicroPhys-BMS achieves $93.83\%$ CDF within $|\Delta\text{SOC}| \le 2.0\%$ on Stanford, outperforming PINN ($85.61\%$) and Transformer ($79.33\%$). On CALCE A123, all models achieve $>97.5\%$. Under $\pm 5.0$~mV noise, the proposed model maintains MAE $= 1.343\%$ (Stanford) and $0.249\%$ (CALCE), both within the $1.456\%$ target. As shown in Fig.~\ref{fig:cdf_comparison}(b), the 961-parameter Micro-MLP is the only model satisfying embedded constraints (Flash $< 16$~kB, ASIL-D) while maintaining competitive accuracy, whereas LightGBM requires $157{,}500$ parameters.

\begin{figure*}[!t]
	\centering
	\begin{subfigure}[b]{0.48\textwidth}
		\centering
		\includegraphics[width=0.85\textwidth]{Figure_7_CDF_Comparison.pdf}
		\caption{CDF of absolute SOC error}
		\label{fig:cdf_comparison_a}
	\end{subfigure}
	\hfill
	\begin{subfigure}[b]{0.48\textwidth}
		\centering
		\includegraphics[width=0.85\textwidth]{Figure_8_Hardware_Tradeoff_Comparison.pdf}
		\caption{Accuracy-complexity trade-off}
		\label{fig:hardware_tradeoff_comparison}
	\end{subfigure}
	\caption{Comparative evaluation. (a) CDFs on both datasets: MicroPhys-BMS achieves $93.83\%$ (Stanford) and $98.89\%$ (CALCE) within $|\Delta\text{SOC}| \le 2.0\%$. (b) Accuracy-complexity: the 961-param Micro-MLP resides in the ideal embedded domain (Flash $< 16$~kB), while baselines need $13{,}185$--$157{,}500$ params.}
	\label{fig:cdf_comparison}
\end{figure*}
